% Original: PDF strana 11, gornji slajd; drugi deo.
\slajdid{03-s021b}
\begin{frame}[t,label=03-s021b]{Visoka impedansa i simboli kodera}
\setlength{\parskip}{0pt}
% element: 03-s021b:e03
Mreža se može dodatno pojednostaviti ako se \(A_1,A_0\) neaktivnog kodera
postavljaju u \textcolor{DEisticanje}{\textbf{stanje visoke impedanse}}.
Tada je moguće direktno povezati odgovarajuće izlaze:
\textbf{četvoroulazni ILI gejtovi nisu potrebni}.
\par\vspace{2mm}
\begin{columns}[c,onlytextwidth]
\begin{column}{.36\textwidth}\centering
\Oznaka{Pravougaoni simbol komponente}
\par\vspace{1mm}
\begin{tikzpicture}[x=1mm,y=1mm,font=\fontsize{9}{11}\selectfont]
\draw[DEvod] (0,0) rectangle (28,32);
\node at (14,23) {KP};\node at (14,18) {4/2};
\foreach \lab/\yy in {I3/26,I2/21,I1/16,I0/11,EI/6}{
 \node[anchor=west,inner sep=1pt] at (0,\yy) {\lab};
 \draw[DEvod] (-3,\yy)--(0,\yy);
}
\foreach \lab/\yy in {GS/26,A1/20,A0/14,EO/6}{
 \node[anchor=east,inner sep=1pt] at (28,\yy) {\lab};
 \draw[DEvod] (28,\yy)--(31,\yy);
}
\end{tikzpicture}

\end{column}
\begin{column}{.60\textwidth}
% element: 03-s021b:e04
\textcolor{DEisticanje}{\textbf{Ne postoji jedan strog standard za raspored MSI simbola.}}
\par\vspace{2mm}
Uobičajeno je da se ulazi crtaju na jednoj, a izlazi na drugoj strani,
bez njihovog mešanja.
\par\vspace{2mm}
Najčešće su \textbf{ulazi levo, a izlazi desno}.
Svi priključci moraju imati jasne oznake da bi se šema mogla pratiti.
\end{column}
\end{columns}
\end{frame}
